\section{Exchangeability and de Finetti representation}
\label{sec:annex-p15-exchangeability-and-definetti}

\begin{annexitem}{annex:p15.exchangeable-canonical}{thm:p15-exchangeable-canonical}{Finite-vector symmetry on the full sequence space}
\begin{proof}
A coordinate permutation is a measurable bijection, since the inverse
image of every coordinate cylinder is another such cylinder.
If finite vectors have the stated symmetry, a finite permutation
preserves probabilities of cylinders: include its support and the
specified cylinder coordinates in one initial finite block and sum
or marginalize the unspecified coordinates.
The laws \(P\) and \(P\circ T_\pi^{-1}\) then agree on the cylinder
pi-system and hence on the product sigma-field.
The converse follows by applying this invariance to cylinders.

Invariant events are closed under complements and countable unions,
so each \(\mathcal E_n\) is a sigma-field.
Invariance under the larger permutation group at \(n+1\) implies
invariance under its subgroup at \(n\), proving decrease.
Every finite permutation is contained in some such initial group,
so their intersection is exactly \(\mathcal E\).
An event in the tail field is determined by coordinates beyond each
fixed \(m\); choosing \(m\) beyond a permutation's support proves
its exact invariance. Thus \(\mathcal T\subseteq\mathcal E\).

For a sequence \(Z\) on another probability space, its coordinate map
is measurable and has the canonical sequence law by definition.
Integrals of functions of the sequence and probabilities of sequence
events therefore agree with their canonical integrals.
For a conditional identity given \(\mathcal E\), integrate on
each \(A\in\mathcal E\) and pull the integral through this map.
The pulled-back candidate is measurable with respect to the
pulled-back sigma-field and has the required all-event identity.
This proves the stated reduction without assuming a regular
conditional probability for every event of the original sample space.
\end{proof}
\end{annexitem}

\begin{annexitem}{annex:p15.exchangeable-averages}{thm:p15-exchangeable-averages}{Exact conditional symmetrization}
\begin{proof}
Permuting the first \(n\) coordinates merely permutes the summands
indexed by \(I_{n,k}\), so \(A_{n,k}(\varphi)\) is
\(\mathcal E_n\)-measurable. Each summand has the same integrable
distribution as \(\varphi(X_1,\ldots,X_k)\), by exchangeability.
Hence the finite average is integrable.

For \(A\in\mathcal E_n\) and any \(i\in I_{n,k}\), choose a
permutation of \(1,\ldots,n\) sending \(j\) to \(i_j\) for
\(1\leq j\leq k\). Measure invariance and exact invariance of \(A\)
give
\[
 E[\mathbf1_A\varphi(X_{i_1},\ldots,X_{i_k})]
                    =E[\mathbf1_A\varphi(X_1,\ldots,X_k)].
\]
Average over the \((n)_k\) index tuples.
The resulting all-event integral identity proves the conditional
expectation formula. Apply \cref{thm:p15-levy-backward} to the
decreasing fields \(\mathcal E_n\), starting at \(n=k\).
Their intersection is \(\mathcal E\), which gives both stated modes
of convergence. No boundedness is needed for this identity beyond
the declared integrability.
\end{proof}
\end{annexitem}

\begin{annexitem}{annex:p15.independent-conditional-projection}{thm:p15-independent-conditional-projection}{The second-moment argument for a conditional mean}
\begin{proof}
Conditional Jensen gives \(EY^2\leq EV^2<\infty\), and \(EY=EV\).
For the bounded truncations \(Y^{(m)}=\max(-m,\min(Y,m))\),
conditional pullout gives
\(E[(V-Y)Y^{(m)}]=0\).
They converge to \(Y\) in \(L^2\).
Cauchy--Schwarz therefore permits passage to
\(E[(V-Y)Y]=0\), or \(EVY=EY^2\).

Replace \(Y\) in products by its \(\mathcal H\)-measurable version.
Independence from \(V\), with the integrability of the product supplied
by Cauchy--Schwarz, gives
\(EVY=(EV)(EY)=(EY)^2\).
Consequently \(E(Y-EY)^2=0\), proving \(Y=EV\) almost surely.
Versions differing on a null set have the same integrals here.
\end{proof}
\end{annexitem}

\begin{annexitem}{annex:p15.reverse-hewitt-savage}{thm:p15-reverse-hewitt-savage}{Deleting a finite block before the independent projection}
\begin{proof}
Let \(|\varphi|\leq M\), and let
\(Y=E(\varphi(X_1,\ldots,X_k)\mid\mathcal E)\).
The symmetrization theorem gives \(A_{n,k}(\varphi)\to Y\)
almost surely.
For fixed \(m\), keep only tuples with all indices greater than \(m\),
but retain the original denominator:
\[
 D_{n,m}=\frac1{(n)_k}
       \sum_{\substack{i\in I_{n,k}\\i_j>m\ (1\leq j\leq k)}}
                         \varphi(X_{i_1},\ldots,X_{i_k}).
\]
This is measurable with respect to \(\sigma(X_j:j>m)\).
A uniform ordered distinct tuple has probability at most \(km/n\)
of meeting the first \(m\) indices, by the union bound over its
\(k\) positions. Thus for \(n\geq m+k\),
\[
                 |A_{n,k}(\varphi)-D_{n,m}|\leq Mkm/n.
\]
The bounded variable \(\limsup_nD_{n,m}\) is a version of \(Y\)
measurable with respect to that tail coordinate field.
No equality of uncompleted sigma-fields is inferred from this
almost-sure version statement.

Take \(m=k\). The first \(k\) i.i.d. coordinates are independent of
the sigma-field of all later coordinates; this follows from finite
independence and the cylinder pi-lambda argument.
The preceding lemma applied to
\(V=\varphi(X_1,\ldots,X_k)\) therefore gives
\(Y=EV\). This proves the bounded-average assertion.

For a finite-coordinate event \(B\), take its indicator as \(\varphi\).
The identity \(E(\mathbf1_B\mid\mathcal E)=P(B)\) says that every
\(A\in\mathcal E\) is independent of \(B\).
For fixed \(A\), the class of \(B\) satisfying
\(P(A\cap B)=P(A)P(B)\) is a lambda-system containing all initial
coordinate cylinders, hence all of \(\mathcal F\).
Take \(B=A\). Then \(P(A)=P(A)^2\), so \(P(A)\) is zero or one.
This proof used neither an earlier zero-one law nor a strong law.
\end{proof}
\end{annexitem}

\begin{annexitem}{annex:p15.reverse-strong-law}{thm:p15-reverse-strong-law}{The partial-sum reverse martingale}
\begin{proof}
Successive differences of the sums recover each \(X_{n+j}\), \(j\geq1\);
conversely those coordinates and \(S_n\) recover all later sums.
This proves the equality defining \(\mathcal H_n\).
It also proves \(\mathcal H_{n+1}\subseteq\mathcal H_n\).
Every generating function is invariant under permutations of the first
\(n\) coordinates, so \(\mathcal H_n\subseteq\mathcal E_n\).

For \(j\leq n\), permutation invariance of every
\(A\in\mathcal H_n\) gives equal integrals of
\(\mathbf1_A X_j\). Hence their conditional expectations given
\(\mathcal H_n\) are equal. Their sum is
\(E(S_n\mid\mathcal H_n)=S_n\), so each equals \(S_n/n\).
In particular the averages are \(E(X_1\mid\mathcal H_n)\).
The downward conditional-limit theorem gives convergence almost surely
and in \(L^1\) to \(E(X_1\mid\bigcap_n\mathcal H_n)\).
The intersection is contained in \(\mathcal E\), which is
probability-trivial by the just-proved Hewitt--Savage theorem.
Its integrable conditional mean is consequently a constant, and its
expectation identifies that constant as \(EX_1\).
This proves the full source reverse-martingale route independently
of the earlier truncation and random-series strong-law arguments.
\end{proof}
\end{annexitem}

\begin{annexitem}{annex:p15.reverse-ballot}{thm:p15-reverse-ballot}{Bounded reverse stopping after conditioning on the final sum}
\begin{proof}
If \(S_N\geq N\), the last inequality defining \(G\) fails.
For any integer \(m<N\) with \(P(S_N=m)>0\), work under the
conditional probability \(Q_m=P(\,\cdot\mid S_N=m)\).
The conditional vector remains exchangeable, because its conditioning
event is symmetric. Its nonnegative coordinates are all bounded by \(m\),
so all reverse averages are integrable under this probability.

For \(0\leq r\leq N-1\), set
\[
       Y_r=S_{N-r}/(N-r),\qquad
       \mathcal H_r=\sigma(S_{N-r},\ldots,S_N).
\]
This filtration increases with \(r\).
For \(k=N-r\geq2\), the information in \(\mathcal H_r\)
specifies \(S_k\) and the later increments and is invariant under
permutations of the first \(k\) coordinates.
Their conditional means are therefore all \(S_k/k\).
It follows that
\[
 E_{Q_m}\!\left(\frac{S_{k-1}}{k-1}\,\middle|\,\mathcal H_r\right)
   =\frac{S_k-S_k/k}{k-1}=\frac{S_k}{k}.
\]
Thus \(Y_r\) is a martingale, with \(Y_0=m/N<1\).

Stop at the first \(r\) with \(Y_r\geq1\), using \(N-1\) if
there is no such \(r\). This is a bounded stopping time \(\tau\).
On \(G\), no threshold is crossed and \(Y_\tau=S_1=0\), since
\(0\leq S_1<1\) is an integer.
On \(G^c\), the first threshold encountered while reading the sums
backwards is at some \(j<N\) with \(S_j\geq j\) and
\(S_{j+1}<j+1\). Integer values and nonnegative increments imply
\(S_j\leq S_{j+1}\leq j\); hence \(S_j=j\) and \(Y_\tau=1\).
Therefore \(Y_\tau=\mathbf1_{G^c}\).
Bounded optional sampling gives \(Q_m(G^c)=m/N\).
The case \(N=1,m=0\) has \(G\) certain and gives the same conclusion.
Together with the \(m\geq N\) cases, this proves the conditional formula
on every positive-probability value of the integer-valued final sum.
There was no use of an unconditional first moment.

For the election, \(S_N=2\beta\) and
\(S_j<j\) means that the number of B votes among the first \(j\)
is strictly less than the number of A votes.
The conditional formula is \(1-2\beta/(\alpha+\beta)\), as stated.
For the real one-draw example the inequality \(1/2<1\) always holds.
For the two signed draws, conditional on zero sum the orders
\((-1,1)\) and \((1,-1)\) are equally likely; only the first satisfies
the strict inequalities. These give both boundary counterexamples.
\end{proof}
\end{annexitem}

\begin{annexitem}{annex:p15.definetti-factorization}{thm:p15-definetti-factorization}{The exact repeated-index remainder}
\begin{proof}
Let \(k\geq2\), with bounded measurable
\(f:S^{k-1}\to\mathbb R\), \(g:S\to\mathbb R\).
Set \(\varphi(x_1,\ldots,x_k)=f(x_1,\ldots,x_{k-1})g(x_k)\),
and \(\varphi_j(x_1,\ldots,x_{k-1})=f(x_1,\ldots,x_{k-1})g(x_j)\).
Multiplying the two finite sums and dividing according to whether the
new index is distinct from the first \(k-1\) gives exactly
\[
 \begin{aligned}
 &(n)_{k-1}n A_{n,k-1}(f)A_{n,1}(g)\\
 &\qquad=(n)_k A_{n,k}(\varphi)
                +(n)_{k-1}\sum_{j=1}^{k-1}A_{n,k-1}(\varphi_j).
 \end{aligned}
\]
For \(n\geq k\), division by \((n)_k=(n)_{k-1}(n-k+1)\) yields
\[
 A_{n,k}(\varphi)
   =\frac{n}{n-k+1}A_{n,k-1}(f)A_{n,1}(g)
       -\frac1{n-k+1}\sum_{j=1}^{k-1}A_{n,k-1}(\varphi_j).
\]
Every remainder average has absolute value at most
\(\|f\|_\infty\|g\|_\infty\), so the last term tends to zero.
The symmetrization theorem gives the almost-sure limits of the
finitely many averages used here, on their common full-measure event.
Taking limits proves
\[
 E(f(X_1,\ldots,X_{k-1})g(X_k)\mid\mathcal E)
   =E(f(X_1,\ldots,X_{k-1})\mid\mathcal E)
                    E(g(X_1)\mid\mathcal E).
\]
For one coordinate, finite permutation invariance on every
\(A\in\mathcal E\) already identifies the conditional laws of \(X_j\)
and \(X_1\). Induct on \(k\), taking \(f=\prod_{j<k}f_j\) and \(g=f_k\),
to obtain the displayed full product identity.
All these are scalar conditional identities with finite specified
families of functions; no regular conditional measure on an arbitrary
target space has been assumed.
\end{proof}
\end{annexitem}


\begin{annexitem}{annex:p15.definetti-kernel}{thm:p15-definetti-kernel}{A measurable product law at every conditioning argument}
\begin{proof}
Regular conditional laws into standard Borel spaces were constructed in
\cref{sec:p15-regular-conditional-laws}, so choose \(K\) as stated.
For every probability \(\mu\) on \(S\), its countable independent
product exists by the Borel coding and extension construction
\cref{thm:p15-coded-space-extension}.
For a cylinder \(D=\{x_1\in B_1,\ldots,x_k\in B_k\}\),
\[
                        K_\omega^{\otimes\mathbb N_+}(D)
                                      =\prod_{j=1}^kK_\omega(B_j)
\]
is \(\mathcal E\)-measurable in \(\omega\).
The class of sequence events \(D\) for which that measurability holds
is a lambda-system: complement gives one minus the value, and disjoint
countable unions give pointwise sums.
It contains the cylinder pi-system and hence every product-measurable
event. Thus \(J\) is a probability kernel, with a probability measure
at every argument.

The conditional-product identity gives, for every \(A\in\mathcal E\)
and every cylinder \(D\),
\[
                          P(A\cap D)=\int_AJ(\omega,D)\,dP(\omega).
\]
For fixed \(A\), both sides are finite measures in \(D\).
They agree on the cylinder pi-system and therefore on all of
\(\mathcal F\). This is precisely the regular conditional-law identity
for the complete sequence given \(\mathcal E\).

The same lambda-system argument, now starting from the measurable
evaluations \(e_B(\mu)\), shows that
\(\mu\mapsto\mu^{\otimes\mathbb N_+}(D)\) is measurable on
\(\mathscr P(S)\) for every \(D\in\mathcal F\).
The kernel \(K\) is a measurable map into that evaluation space
by its definition. Taking its distribution \(\Lambda\), the case
\(A=\Omega\) of the conditional identity and the pushforward
integration formula give the asserted mixture on all sequence events.
For a cylinder this is exactly the displayed product integral.
\end{proof}
\end{annexitem}

\begin{annexitem}{annex:p15.directing-law-uniqueness}{thm:p15-directing-law-uniqueness}{Countable recovery and equality of information modulo null events}
\begin{proof}
A countable determining pi-system \(\mathcal C\) exists by pulling
rational half-lines back through a Borel real coding of \(S\) and
including \(S,\varnothing\).
The sigma-field of evaluations on \(\mathscr P(S)\) is already
generated by \(e_C\), \(C\in\mathcal C\).
Indeed the sets \(B\in\mathcal S\) whose evaluation \(e_B\) is measurable
with respect to these countably many evaluations form a lambda-system,
using \(e_{S\setminus B}=1-e_B\) and sums on disjoint unions.
The pi-lambda theorem proves the assertion for every \(B\).

Apply the symmetrization theorem with \(k=1,\varphi=\mathbf1_C\).
Its limit is \(K(C)\) almost surely. Intersect over the countable
class to obtain simultaneous convergence.
For every fixed \(m\), deletion of the first \(m\) terms changes an
empirical frequency by at most \(m/n\).
Consequently its limsup, defined on every sequence, is measurable with
respect to \(\sigma(X_j:j>m)\), and hence with respect to \(\mathcal T\).
It is a tail-measurable version of \(K(C)\).

The class \(\mathcal G^P\) in the definition is a sigma-field:
complements preserve the symmetric-difference bound, and a countable
union of chosen representatives differs only on the union of their
null differences. It contains every ambient measurable null event.
Thus the preceding versions and the countable evaluation generators
give \(\sigma(K)\subseteq\mathcal T^P\), hence
\(\sigma(K)^P\subseteq\mathcal T^P\).
The exact inclusion \(\mathcal T\subseteq\mathcal E\) gives the next
inclusion into \(\mathcal E^P\).

For the reverse inclusion take \(A\in\mathcal E\).
For every fixed \(\mu\in\mathscr P(S)\), the coordinates under
\(\mu^{\otimes\mathbb N_+}\) are i.i.d., so the source Hewitt--Savage
proof gives
\[
                            h_A(\mu)=\mu^{\otimes\mathbb N_+}(A)
                                                   \in\{0,1\}.
\]
The function is measurable by the product-kernel proof.
As \(J\) is conditional given \(\mathcal E\),
\(h_A(K)=J(\cdot,A)=\mathbf1_A\), \(P\)-almost surely.
The set \(\{h_A(K)=1\}\in\sigma(K)\) is therefore a representative
of \(A\) modulo a null event. This proves the three-way equality.
No literal equality of the unaugmented fields was used.

To prove uniqueness of the mixing law, let another probability
\(\Lambda'\) give the same finite-coordinate mixture probabilities.
The product-kernel theorem constructs the joint probability
\(\Lambda'(d\mu)\mu^{\otimes\mathbb N_+}(dx)\)
by \cref{thm:p15-kernel-joint-measure}.
Its sequence marginal is \(P\), by uniqueness from cylinders.
For every fixed \(\mu\), the i.i.d. strong law for each indicator
\(\mathbf1_C\), simultaneously over the countable class, gives
limiting frequencies \(e_C(\mu)\).
This convergence event is measurable in \((\mu,x)\), using countably
many coordinates and the evaluations. Kernel integration therefore
gives joint probability one to it.
The distribution of the vector \((e_C(\mu))_{C\in\mathcal C}\)
under \(\Lambda'\) is consequently the distribution of the frequency
limit vector under \(P\), which is also its distribution under
\(\Lambda\). Since these evaluations generate the sigma-field of
\(\mathscr P(S)\), the two mixing measures agree.

Finally, if \(\Lambda\) is concentrated at one law \(\mu_0\),
the mixture formula says \(P=\mu_0^{\otimes\mathbb N_+}\).
Conversely under an i.i.d. law, \(\mathcal E\) is probability-trivial.
Each real variable \(K(C)\) is then almost surely constant, with value
\(E K(C)=P(X_1\in C)\).
Countability and the determining property give \(K=\mathcal L(X_1)\)
almost surely, so \(\Lambda\) is a point mass.
The same last argument shows that probability-triviality of
\(\mathcal E\) for an exchangeable law implies that conclusion.
\end{proof}
\end{annexitem}

\begin{annexitem}{annex:p15.exchangeable-strong-laws}{thm:p15-exchangeable-strong-laws}{Integrable averages and empirical weak convergence}
\begin{proof}
For the first assertion use the symmetrization theorem with \(k=1\)
and \(\varphi=f\). Its average is precisely the sample average,
and it converges almost surely and in \(L^1\).
The conditional-kernel integration theorem identifies the limit with
\(\int f\,dK\) on the full-measure set where the absolute integral
is finite.

For the weaker, conditional-integrability hypothesis let
\(B=\{\omega:\int|f|\,dK_\omega<\infty\}\in\mathcal E\).
It has probability one by the law \(\Lambda\) of \(K\).
Define \(m_f(\omega)=\int f\,dK_\omega\) on \(B\), and zero elsewhere.
It is a finite measurable version.
For each \(\omega\in B\), the ordinary i.i.d. strong law under
\(J(\omega,\cdot)=K_\omega^{\otimes\mathbb N_+}\) gives convergence
of the sample averages to \(m_f(\omega)\) with probability one.
The event of that convergence is measurable in
\((\omega,x)\in(\Omega,\mathcal E)\times S^{\mathbb N_+}\):
write the usual countable limit inequalities with center \(m_f(\omega)\).
The joint law of the conditioning argument and its actual sequence
is \(P(d\omega)J(\omega,dx)\), as follows first on rectangles from the
conditional identity and then by pi-lambda uniqueness.
Integrating the section probabilities proves the almost-sure
assertion for the actual sequence. No unconditional first moment
or \(L^1\) conclusion is inferred in this case.

For \(S=\mathbb R\), choose rational half-lines in the determining
class. Outside one null event, the empirical CDFs
\(\widehat F_n(q)\) converge to \(F_\omega(q)=K_\omega((-\infty,q])\)
for every rational \(q\).
Fix any continuity point \(x\) of this limiting CDF.
For rational \(q<x<r\),
\(\widehat F_n(q)\leq\widehat F_n(x)\leq\widehat F_n(r)\).
Take liminf and limsup in \(n\), then let rationals \(q\) increase
to \(x\) and \(r\) decrease to \(x\).
Continuity at \(x\) and right continuity of a CDF make both bounds
equal \(F_\omega(x)\).
This proves CDF convergence at every continuity point on that same
event, and therefore weak convergence of the empirical laws.
\end{proof}
\end{annexitem}

\begin{annexitem}{annex:p15.binary-definetti}{thm:p15-binary-definetti}{The binary parameter and uniform count orbits}
\begin{proof}
For \(S=\{0,1\}\), every probability measure is determined by its
success probability. Put \(\Theta=K(\{1\})\).
The conditional-product formula gives a word of count \(s\)
conditional probability \(\Theta^s(1-\Theta)^{n-s}\).
Its expectation is the displayed integral for
\(\lambda=\mathcal L(\Theta)\).
This convention gives the correct deterministic Bernoulli law at
either endpoint. The empirical-frequency theorem gives the
almost-sure limit \(\Theta\), so its law identifies \(\lambda\)
uniquely; it is also the special case of directing-law uniqueness.
There are \(\binom ns\) disjoint words with that count, proving
the count formula.

In the finite vector, permutations act transitively on the
\(\binom Nm\) words with count \(m\). Their unconditional
probabilities are equal by exchangeability.
When the count event has positive probability, division by its
probability makes each conditional word probability \(1/\binom Nm\).
Exactly \(\binom{N-k}{m-k}\) of those words start with \(k\) ones;
if \(m<k\) none does. Counting proves the conditional formula.
The argument does not condition at a count value of probability zero
and does not assume that a finite vector has an infinite extension.
\end{proof}
\end{annexitem}

\begin{annexitem}{annex:p15.finite-urn-approximation}{thm:p15-finite-urn-approximation}{Coupled indices and compact selection of the mixing law}
\begin{proof}
Fix a binary population \(z_1,\ldots,z_N\) with \(m\) ones.
Let \(I_1,\ldots,I_k\) be independent uniform indices in
\(\{1,\ldots,N\}\). On the event \(D\) that they are all distinct,
their ordered tuple is uniform over the injective tuples.
Define an injective tuple \(J\) to equal \(I\) on \(D\);
on \(D^c\), draw \(J\) independently and uniformly among injective tuples.
Its unconditional law is uniform among those tuples, because both
conditional constructions have that same law.
Thus \((z_{J_1},\ldots,z_{J_k})\) has the without-replacement law
and \((z_{I_1},\ldots,z_{I_k})\) has independent Bernoulli law
of probability \(m/N\). They agree on \(D\).
For every event the difference of its probabilities is at most
\(P(D^c)\), and the union bound gives
\[
                    P(D^c)
                       \leq\sum_{i<j}P(I_i=I_j)
                       =\frac{k(k-1)}{2N}.
\]
This proves the event-supremum total variation bound, including the
all-zero and all-one populations.

For a finite exchangeable vector, condition on its final count.
The preceding uniform-orbit calculation identifies the conditional
first-\(k\) law as \(H_{m,N}^{(k)}\).
Average the event inequality over the count distribution; the bound
does not depend on \(m\). This proves the mixed bound.

For an infinite exchangeable binary sequence let
\(\lambda_N=\mathcal L(S_N/N)\).
All these laws are supported on the compact interval \([0,1]\),
hence are tight. The already proved real-law selection theorem
gives a subsequence converging weakly to a probability \(\lambda\)
supported there. The support follows, for example, by the closed-set
bound for \([0,1]\).
Fix any word of length \(k\) with count \(s\).
The mixed bound says that its fixed sequence probability differs from
\(\int\theta^s(1-\theta)^{k-s}\,d\lambda_N(\theta)\)
by at most \(k(k-1)/(2N)\).
The integrand is a continuous polynomial on \([0,1]\), including
the stated endpoint convention. Weak convergence along the selected
subsequence therefore gives the binary mixture formula for this word.
The same selected \(\lambda\) works for every fixed finite word.

This is an existence proof independent of the general conditional-product
theorem. Its mixing law can also be identified without that theorem:
choose a parameter with law \(\lambda\) independently of a uniform
sequence and threshold the uniforms at that parameter.
Its finite word probabilities are exactly the mixture integrals,
so its full sequence law is the original law by cylinder uniqueness.
The measurable-section strong-law argument used for the Pólya mixture
then shows that its frequency limit equals the parameter almost surely.
The original sequence therefore has a limiting frequency with law
\(\lambda\), proving uniqueness for this construction too.
\end{proof}
\end{annexitem}

\begin{annexitem}{annex:p15.exchangeable-second-moments}{thm:p15-exchangeable-second-moments}{Truncation in the product identity and the integrable pair average}
\begin{proof}
Let \(g_b(x)=\max(-b,\min(x,b))\), and
\(M_b=E(g_b(X_1)\mid\mathcal E)\).
The bounded conditional-product identity gives
\(E[g_b(X_1)g_b(X_2)]=EM_b^2\).
As \(b\to\infty\), the truncations converge to \(X_j\) in \(L^2\).
Conditional Jensen gives
\(E|M_b-M|^2\leq E|g_b(X_1)-X_1|^2\to0\).
Cauchy--Schwarz then gives \(M_b^2\to M^2\) in \(L^1\).
It also gives convergence of the products on the left in \(L^1\):
split their difference into
\((g_b(X_1)-X_1)g_b(X_2)+X_1(g_b(X_2)-X_2)\)
and bound each term by its two \(L^2\) norms.
Taking limits proves \(E(X_1X_2)=EM^2\).
Since \(EM=EX_1=EX_2\), subtracting the square of the common mean
gives the covariance identity and its nonnegativity.

For the i.i.d. pair statistic, \(\varphi(x,y)=(x-y)^2\) is integrable,
because it is bounded by \(2x^2+2y^2\) at the coordinate variables.
Its ordered distinct-index average is exactly
\(\binom n2^{-1}\sum_{i<j}(X_i-X_j)^2\).
The integrable symmetrization theorem gives almost-sure and \(L^1\)
convergence to \(E(\varphi(X_1,X_2)\mid\mathcal E)\).
Hewitt--Savage makes this integrable conditional mean constant.
Independence gives
\(E(X_1-X_2)^2=2EX_1^2-2(EX_1)^2=2\sigma^2\).

For the stated finite pair, each coordinate has mean \(1/2\),
while their product is always zero, so the covariance is \(-1/4\).
An infinite exchangeable extension would contradict the nonnegative
covariance just proved. A Bernoulli mixture would have
\(P(X_1=X_2=1)=E\Theta^2=0\), forcing \(\Theta=0\) almost surely
and contradicting its marginal success probability \(1/2\).
This proves both obstructions directly.
\end{proof}
\end{annexitem}

\begin{annexitem}{annex:p15.backward-lp-limit}{thm:p15-backward-lp-limit}{Truncation for finite powers and the bounded counterexample}
\begin{proof}
The downward conditional theorem supplies the almost-sure and \(L^1\)
limit. For \(p>1\), let \(Z^{(b)}=\max(-b,\min(Z,b))\),
\(Y_n^{(b)}=E(Z^{(b)}\mid\mathcal G_n)\) and
\(Y_\infty^{(b)}=E(Z^{(b)}\mid\bigcap_n\mathcal G_n)\).
These bounded conditional variables converge almost surely, hence
in \(L^p\) by domination by \((2b)^p\).
Conditional Jensen gives
\[
 E|Y_n-Y_n^{(b)}|^p\leq E|Z-Z^{(b)}|^p,
 \qquad
 E|Y_\infty-Y_\infty^{(b)}|^p\leq E|Z-Z^{(b)}|^p .
\]
Use \(|u+v+w|^p\leq3^{p-1}(|u|^p+|v|^p+|w|^p)\).
For fixed \(b\), the middle bounded difference tends to zero, so
\[
 \limsup_n E|Y_n-Y_\infty|^p
                     \leq2\cdot3^{p-1}E|Z-Z^{(b)}|^p.
\]
The right side tends to zero by integrability of \(|Z|^p\).
For \(0<p\leq1\), concavity gives
\(E|Y_n-Y_\infty|^p\leq(E|Y_n-Y_\infty|)^p\to0\).

For the example, \(\mathcal G_n\) decreases.
Every point of \((a_n,1]\) belongs to the same atom of
\(\mathcal G_n\), since all its generating intervals lie in
\((0,a_n]\). On the latter set, \(Z=1\); on the former atom
its conditional average is
\((1/2-a_n)/(1-a_n)\). This proves the displayed \(Y_n\)
by its all-event integral identity.
Any two positive points lie together in the outer atom for all
sufficiently large \(n\). Hence an event in every \(\mathcal G_n\)
either contains both or neither, and the intersection is the trivial
sigma-field. Its conditional mean is \(EZ=1/2\).
The difference equals \(1/2\) on the positive-measure interval
\((0,a_n]\). On its complement the absolute difference is
\(a_n/(2(1-a_n))\leq1/2\).
Thus the essential supremum is exactly \(1/2\) at every time.
\end{proof}
\end{annexitem}

\begin{annexitem}{annex:p15.backward-changing-variables}{thm:p15-backward-changing-variables}{A decreasing oscillation envelope}
\begin{proof}
The almost-sure limit satisfies \(|Z|\leq W\).
For \(N\geq0\), set
\(R_N=\sup_{n\geq N}|Z_n-Z|\).
These measurable nonnegative variables decrease to zero almost surely
and are bounded by \(2W\), so \(ER_N\to0\).
For \(n\geq N\), conditional monotonicity and the absolute-value
inequality give
\[
 |E(Z_n\mid\mathcal G_n)-E(Z\mid\mathcal G_n)|
                              \leq E(R_N\mid\mathcal G_n).
\]
All pairs \(n,N\) are countable, so their exceptional sets can be
removed together.
For each fixed \(N\), the downward conditional theorem makes the
right side converge almost surely to \(E(R_N\mid\mathcal G_\infty)\).
These latter conditional variables decrease to zero almost surely:
their decreasing limit is nonnegative and has expectation at most
\(\lim_NER_N=0\).
Taking limsup in \(n\) and then the infimum over \(N\) therefore
makes the displayed difference tend to zero.
The downward theorem applied to the fixed \(Z\) identifies the
remaining limit as \(E(Z\mid\mathcal G_\infty)\).

Finally conditional \(L^1\) contraction gives
\[
 E|E(Z_n\mid\mathcal G_n)-E(Z\mid\mathcal G_n)|
                                      \leq E|Z_n-Z|\longrightarrow0.
\]
Combine this with the fixed-variable \(L^1\) downward limit to
obtain the stated \(L^1\) convergence.
\end{proof}
\end{annexitem}

